\subsection{The phase-label theorem: program elimination and flat components}
\label{ian:phase-label-theorem}

Status: proved. Theorems~\ref{ian:phase-label-elimination} and \ref{ian:phase-label-hard-sector-converse} and Proposition~\ref{ian:phase-label-flat-components} hold for every decomposition satisfying \eqref{ian:phase-label-component-bound}; every application with product or exchangeable components (the noisy, soft, bounded-degree local and positive instances) uses the uniform-input simulator \(s_U\) of Subsection~\ref{ian:phase-label-admissibility}, which Corollary~\ref{cor:resolution-product-alphabet-aa} supplies, and so rests only on inequality (BC) of Huffstutler--Miloschewsky--Podder, the sole import of that corollary; no influence form of uniform AA is used. The classical lower bounds \eqref{ian:phase-label-clean-bound}, \eqref{ian:phase-label-noisy-bound} and \eqref{ian:phase-label-local-bound} are unconditional. Failure of full AA is equivalent to hard sectors and implies a hidden (label, transcript) explanation on a hard sector (the explanation alone does not imply failure); automatic label-only explanation is refuted on a prescribed hard sector, and no converse in terms of natural phases is proved.

Fix a family \((\mu_N)\) and a decomposition satisfying
\eqref{ian:phase-label-component-bound} with polynomial \(\Psi\). By
Lemma~\ref{ian:phase-label-measurable-choice}, after replacing
\(\Psi(T,y)\) by \(\Psi(T,2y)\) for uncountable label spaces, there are
measurable component simulators
\[
 D_z\in\mathcal D_s,\qquad
 s=\lceil\Psi(T,\delta^{-1})\rceil,\qquad
 \E_{\mu_z}|D_z-Q|\le\delta .
\]
On the joint space \(\Pr(Z\in dz,X=x)=\pi(dz)\mu_z(x)\) put
\[
 \widehat Q=D_Z(X),
 \tag{20.1}\label{ian:phase-label-program}
\]
so \(\E|\widehat Q-Q|\le\delta\). This is a program indexed by \(Z\) and
applied to \(X\), not a function of \(Z\) alone. Note that
\[
 \E[Q(X)\mid Z,X]=Q(X),
 \tag{20.2}\label{ian:phase-label-conditioning-correction}
\]
so conditioning on \((Z,X)\) attributes no variance to the label; the
label-only quantity is \(\E[Q\mid Z]\), or \(\E[Q\mid Z,H]\) on a sector.

\begin{theorem}[Eliminating a component program]
\label{ian:phase-label-elimination}
If a depth-\(t_0\) tree produces a surrogate label \(\widehat Z\) with
\(\E|D_{\widehat Z}(X)-D_Z(X)|\le\zeta\) (for instance
\(\Pr(\widehat Z\ne Z)\le\zeta\)), then a tree of depth at most
\(t_0+s\) approximates \(Q\) on \(\mu\) with error at most
\(\delta+\zeta\).

With \(L_{\widehat Q}(t)=\inf_{D\in\mathcal D_t}\E|D(X)-\widehat Q|\),
\[
 \bigl|L_{\widehat Q}(t)-\ell_\mu(t,Q)\bigr|\le\delta
 \qquad(t\ge0).
 \tag{20.3}\label{ian:phase-label-loss-comparison}
\]
Consequently full AA holds if and only if there is a polynomial \(\Phi\),
fixed before \(N,T,\varepsilon,Q\), such that for all \(N,T\ge1\),
\(\varepsilon\in(0,1)\) and \(Q\in\mathcal Q_{N,T}\), with
\(\widehat Q\) formed at \(\delta=\varepsilon/4\), some
\(G\in\mathcal D_{\lceil\Phi(T,\varepsilon^{-1})\rceil}\) satisfies
\[
 \E|G(X)-\widehat Q|\le\varepsilon/2.
 \tag{20.4}\label{ian:phase-label-program-eliminability}
\]
This holds for every fixed measurable choice of component simulators.
\end{theorem}

\begin{proof}
Run the surrogate tree, then \(D_{\widehat Z}\); the triangle inequality
gives the error. For each \(D\) the two losses in
\eqref{ian:phase-label-loss-comparison} differ by at most
\(\E|\widehat Q-Q|\le\delta\). If full AA holds with polynomial
\(\Phi_0\), its simulator at accuracy \(\varepsilon/4\) has depth
\(\lceil\Phi_0(T,4\varepsilon^{-1})\rceil\) and distance at most
\(\varepsilon/2\) to \(\widehat Q\); take \(\Phi(T,y)=\Phi_0(T,4y)\).
Conversely \eqref{ian:phase-label-program-eliminability} gives error at
most \(3\varepsilon/4\) for \(Q\).
\end{proof}

Randomized selectors obey the same statements after their
input-independent coins are fixed by averaging. A selector need not
recover the label, only a program compatible with the target.

\begin{theorem}[Explicit hard-sector converse]
\label{ian:phase-label-hard-sector-converse}
Suppose full AA fails for \((\mu_N)\). For every \(\theta_1,\theta_2\ge1\)
there exist \(N,T\ge1\), \(\varepsilon\in(0,1/8)\),
\(Q\in\mathcal Q_{N,T}\) and \(H\subseteq\Omega_N\) such that, with
\[
 w_0=\frac{\varepsilon^3}{64},\qquad
 a=\frac{\varepsilon^2}{4},\qquad
 \delta=\frac{\varepsilon^7}{2^{22}},\qquad
 s=\lceil\Psi(T,\delta^{-1})\rceil,\qquad
 t=\left\lceil\theta_1(T+\varepsilon^{-1})^{\theta_2}\right\rceil+s+1,
 \tag{20.5}\label{ian:phase-label-explicit-parameters}
\]
one has
\[
 \mu_N(H)\ge w_0,\qquad
 \operatorname{Hard}(Q,\mu_N(\cdot\mid H);t,a,a/4).
 \tag{20.6}\label{ian:phase-label-hard-witness}
\]
Lift the joint law to the event \(\{X\in H\}\), write \(\E_H\) for its
expectation, and form \(\widehat Q\) at accuracy \(\delta\). With
\(u=a^2/4096=\varepsilon^4/2^{16}\),
\[
 \E_H|Q-\widehat Q|\le u,\qquad
 \Var_H\widehat Q\ge\Var_HQ-2u,
 \tag{20.7}\label{ian:phase-label-lifted-signal}
\]
and every \(F\in\mathcal F_t\) on the physical spins satisfies
\[
 |\Cov_H(\widehat Q,F)|\le a/4+2u.
 \tag{20.8}\label{ian:phase-label-lifted-covariance}
\]
Let \(\tau_Z(X)\) be the transcript of \(D_Z\) on \(X\), and put
\(\widetilde Z=(Z,\tau_Z(X))\), \(M=\E_H[Q\mid\widetilde Z]\). Then
\[
 \E_H\Var(Q\mid\widetilde Z)\le u,\qquad
 \Var_HM\ge\Var_HQ-u,
 \tag{20.9}\label{ian:phase-label-variance-explanation}
\]
\[
 |\Cov_H(M,F)|\le a/4+a/64
 \qquad(F\in\mathcal F_t).
 \tag{20.10}\label{ian:phase-label-explained-signal-hidden}
\]
Conversely, if witnesses \eqref{ian:phase-label-hard-witness} exist for
all \(\theta_1,\theta_2\ge1\), full AA fails.
\end{theorem}

\begin{proof}
Write \(\Psi(T,y)\le c_\Psi(T+y)^{d_\Psi}\) for \(T,y\ge1\), where
\(c_\Psi\) is the sum of the absolute coefficients and
\(d_\Psi=\max\{1,\deg\Psi\}\). With \(x=T+\varepsilon^{-1}\ge1\),
\(T+\delta^{-1}\le2^{22}x^7\), so
\(t\le\theta_1'x^{\theta_2'}\) for
\(\theta_1'=\theta_1+c_\Psi2^{22d_\Psi}+3\) and
\(\theta_2'=\max\{\theta_2,7d_\Psi\}\). Apply
Theorem~\ref{ian:hard-sector-aa-equivalence}, (1)\(\Rightarrow\)(2), with
budget constants \(\theta_1',\theta_2'\); its proof contains the
reduction to \(\varepsilon<1/8\). With \(\varepsilon=r^{-1}\) it gives
\eqref{ian:phase-label-hard-witness} at depth
\(\lceil\theta_1'x^{\theta_2'}\rceil+1\ge t\), hence at depth \(t\),
since \(\mathcal F_t\subseteq\mathcal F_{t'}\) for \(t\le t'\).

Restriction costs the reciprocal sector weight:
\(\E_H|Q-\widehat Q|\le\delta/w_0=u\). For \(U,V\in[0,1]\) and
\(|F|\le1\) on one probability space,
\(|\Var U-\Var V|\le2\E|U-V|\), from
\(\Var U=\min_{c\in[0,1]}\E(U-c)^2\), and
\(|\Cov(U,F)-\Cov(V,F)|\le2\E|U-V|\). These give
\eqref{ian:phase-label-lifted-signal} and
\eqref{ian:phase-label-lifted-covariance}.

The output \(\widehat Q\) is \(\widetilde Z\)-measurable, so
\(\E_H(Q-M)^2\le\E_H(Q-\widehat Q)^2\le\E_H|Q-\widehat Q|\le u\); total
variance gives \eqref{ian:phase-label-variance-explanation}. Since
\(\E_HM=\E_HQ\),
\(|\Cov_H(M,F)-\Cov_H(Q,F)|=|\E_H[(M-Q)F]|\le\sqrt u=a/64\).

For the converse, \(t\ge\lceil\theta_1(T+\varepsilon^{-1})^{\theta_2}\rceil+1\),
so the witnesses give item (2) of
Theorem~\ref{ian:hard-sector-aa-equivalence} with \(r=\varepsilon^{-1}\),
after the same depth monotonicity.
\end{proof}

\begin{remark}[What is and is not an equivalence]
\label{ian:phase-label-headline}
Equation~\eqref{ian:phase-label-program-eliminability} is an equivalence,
but an operational one: it restates full AA for the fixed programs
\(\widehat Q\). The structural statements are: failure of full AA is
equivalent to polynomially weighted hard sectors
(Theorem~\ref{ian:hard-sector-aa-equivalence}); and failure implies that
on such a sector the label together with its shallow component
transcript, \(\widetilde Z\), explains almost all acceptance variance
while staying invisible to shallow tests
(Theorem~\ref{ian:phase-label-hard-sector-converse}). The explanation
adds nothing in the reverse direction: what forces failure is that the
sector is hard (dispersion plus invisibility), while a
\(\widetilde Z\)-explanation by itself is always available, e.g.\
\(\widetilde Z=X\).
Proposition~\ref{ian:phase-label-flat-components} gives a sufficient
condition for a label-only explanation, and
Proposition~\ref{ian:phase-label-simplex-counterexample} shows it is not
automatic. No claim is made that \(\mu_z(\cdot\mid H)\in\mathcal C\): a
hard sector is an arbitrary event, and only the original component
error and domination by \(1/\mu(H)\) are used.
\end{remark}

\begin{proposition}[A sufficient condition for a label-only statement]
\label{ian:phase-label-flat-components}
Let \(g(Z)=\E[Q\mid Z]\) and \(v=\E\Var(Q\mid Z)=\E(Q-g(Z))^2\). For a
sector of mass \(w=\mu(H)>0\) put \(g_H(Z)=\E[Q\mid Z,H]\). Then
\[
 \E_H\Var(Q\mid Z,H)\le v/w,
 \qquad
 \Var_Hg_H\ge\Var_HQ-v/w.
 \tag{20.11}\label{ian:phase-label-flat-variance}
\]
If \(\operatorname{Hard}(Q,\mu(\cdot\mid H);t,a,\gamma)\), then
\[
 |\Cov_H(g_H,F)|\le\gamma+\sqrt{v/w}
 \qquad(F\in\mathcal F_t).
 \tag{20.12}\label{ian:phase-label-flat-covariance}
\]
In particular \(v\le wa^2/4096\) gives a label-only version of
\eqref{ian:phase-label-variance-explanation}--%
\eqref{ian:phase-label-explained-signal-hidden}. More generally, for any
\(g(Z)\in[0,1]\) with \(e=\E|Q-g(Z)|\),
\[
 \left|
 \inf_{D\in\mathcal D_t}\E|D(X)-g(Z)|
 -\ell_\mu(t,Q)
 \right|\le e,
 \tag{20.13}\label{ian:phase-label-pure-target-comparison}
\]
so a classical label-target lower bound \(e+\Delta\) gives an
acceptance-simulation lower bound \(\Delta\).
\end{proposition}

\begin{proof}
On \(H\) the predictor \(g(Z)\) has squared loss at most \(v/w\), and
\(g_H\) has no larger loss. Total variance and Cauchy--Schwarz, as in the
preceding proof, give the first two claims; the last is the triangle
inequality.
\end{proof}

Component simulability does not give \(v\) small: a shallow component
tree may read an essential random input bit.

In the instances below, \(m\ge2\) is even, \(n=2^m\), \(h\) is the bent
quadratic and \(Q\) is the one-query acceptance function, called \(P\)
in Section~\ref{ian:hard-sector-instances}; put
\(\bar h=\E h(A)=(1-n^{-1/2})/2\) for uniform \(A\).

\begin{proposition}[A hard sector with negligible between-label variance]
\label{ian:phase-label-simplex-counterexample}
Take \(\mu=U_{C_n}\), so \(Q(c_A)=h(A)\), with the
\(\mathcal U\)-decomposition \eqref{ian:phase-label-simplex} and label
\(\xi\). Then
\[
 \Var Q=\frac{1-1/n}{4},
 \qquad
 \Var\bigl(\E[Q\mid\xi]\bigr)
 =\frac14\left(\frac1{m+1}-\frac1n\right).
 \tag{20.14}\label{ian:phase-label-small-between}
\]
Every component computes \(Q\) exactly with one query, yet the fraction
of variance explained by \(\xi\) tends to zero, while \(H=C_n\) is a
weight-one hard sector: \(\operatorname{Hard}(Q,U_{C_n};t,3/16,1/8)\) for
\(n\ge16\), \(t\le m/2-2\).
\end{proposition}

\begin{proof}
Write \(f=(-1)^h\). Flipping coordinate \(i\) of \(A\) multiplies \(f\)
by \((-1)^{A_{i'}}\), \(i'\) the paired coordinate, so
\[
 \E[f(A)\mid\xi]
 =\frac{f(\xi)}{m+1}\Bigl(1+\sum_{i=1}^m(-1)^{\xi_i}\Bigr),
 \qquad
 \E\bigl(\E[f(A)\mid\xi]\bigr)^2=\frac1{m+1},
 \qquad
 \E f(A)=n^{-1/2},
\]
which gives \eqref{ian:phase-label-small-between} via \(Q=(1-f)/2\).
Coordinate universality gives the one-query component simulator, and
hardness is Proposition~\ref{ian:hard-sector-character-instance}, with
covariance at most \((2^t-1)/(2\sqrt n)\) by
Lemma~\ref{ian:hard-sector-orbit-covariance}.
\end{proof}

This refutes automatic between-label concentration on a supplied hard
sector, even for a hereditary decomposition with growing component
min-entropy. It does not exclude some other sector with larger
between-label variance; hard-sector extraction supplies no such sector.

\begin{proposition}[Character and noisy-character obstructions]
\label{ian:phase-label-character-noisy-instances}
For \(U_{C_n}\), one quantum phase query followed by a Walsh transform
returns \(A\), and accepting when \(h(A)=1\) gives \(Q\). Every classical
depth-\(t\) predictor \(D\) of \(h(A)\) satisfies
\[
 \E|D-h(A)|\ge\frac{1-2^{t-m/2}}2.
 \tag{20.15}\label{ian:phase-label-clean-bound}
\]
For \((1-w_n)U_n+w_nU_{C_n}\) the sector \(C_n\) has mass
\(w_n+(1-w_n)n2^{-n}\), at least \(1/2\) in the Gibbs realization at
\(\beta\ge1\) (Proposition~\ref{ian:hard-sector-character-instance}).

For the noisy law \(X=(-1)^Bc_A\eta\), \(\E\eta_z=\rho\in(0,1)\), put
\(g_\rho(A)=\bar h+\rho^2(h(A)-\bar h)\). Then
\[
 \E[Q\mid A,B]=g_\rho(A),\qquad
 \E(Q-g_\rho(A))^2=\frac{2\bar h(1-\bar h)(1-\rho^4)}n,\qquad
 \Var g_\rho=\rho^4\bar h(1-\bar h).
 \tag{20.16}\label{ian:phase-label-noisy-flatness}
\]
Every depth-\(t\) tree satisfies
\(\E|D-g_\rho(A)|\ge\frac{\rho^2}{2}(1-2^{t-m/2})\), hence
\[
 \E|D-Q|
 \ge\frac{\rho^2}{2}(1-2^{t-m/2})
 -\sqrt{\frac{2\bar h(1-\bar h)(1-\rho^4)}n}.
 \tag{20.17}\label{ian:phase-label-noisy-bound}
\]
\end{proposition}

\begin{proof}
After the noise and global sign are granted, a depth-\(t\) transcript
imposes at most \(t\) linear equations on uniform \(A\), i.e.\ restricts
it to an affine subspace of codimension \(c\le t\). There
\(|\E(-1)^{h(A)}|\le2^{c-m/2}\), since all Walsh coefficients of
\((-1)^h\) have magnitude \(2^{-m/2}\); the best leaf output then has
error at least \((1-2^{t-m/2})/2\). The noisy target takes two values
\(\rho^2\) apart, giving its bound.

By the proof of Lemma~\ref{ian:hard-sector-orbit-covariance},
\(Q(\sigma c_A\eta)=\sum_bp_bh(A+b)\) with \(p_b=\widehat\eta(b)^2\).
For independent \(\eta_z\) of mean \(\rho\),
\(\E p_0=\rho^2+(1-\rho^2)/n\) and \(\E p_b=(1-\rho^2)/n\) for \(b\ne0\);
since \(\sum_bh(A+b)=n\bar h\), this gives \(\E[Q\mid A,B]=g_\rho(A)\).
The repo formula \(\Var Q=\bar h(1-\bar h)[\rho^4+2(1-\rho^4)/n]\) (proof of
Proposition~\ref{ian:hard-sector-soft-instance}, from
Theorem~\ref{threshold:noisy-query-theorem}) and total variance give the
rest of \eqref{ian:phase-label-noisy-flatness}; Cauchy--Schwarz gives
\eqref{ian:phase-label-noisy-bound}.
\end{proof}

For noisy inputs the Walsh measurement returns \(A\) with probability
\(\E p_0=\rho^2+(1-\rho^2)/n\); the obstruction concerns the acceptance
probability, not exact recovery. The hard target \(h(A)\) is nonlinear:
every parity \(u\cdot A\) is the single clean coordinate \(c_A(u)\), so
the advantage is not recovery of a classically hidden linear parity.

\begin{proposition}[The soft-code phase label]
\label{ian:phase-label-soft-instance}
Use the joint law \eqref{ian:phase-label-soft-joint}, with
\(V=(-1)^Bc_A\), at fixed \(\beta,L>0\), \(\rho=\tanh(\beta L)\), and
\(\tau_n\) as in Proposition~\ref{ian:hard-sector-soft-instance}. Then
\[
 \E_\mu(Q-g_\rho(A))^2
 \le \frac{2\bar h(1-\bar h)(1-\rho^4)}n+\tau_n,
 \tag{20.18}\label{ian:phase-label-soft-flatness}
\]
and \eqref{ian:phase-label-noisy-bound} holds on \(\mu\) with an
additional loss \(\tau_n\). Using \(s_U\), the components are
uniformly simulable (item 3 after
Lemma~\ref{ian:phase-label-close-products}).
\end{proposition}

\begin{proof}
The joint law \eqref{ian:phase-label-soft-joint} has density
\(\omega(X)\) relative to the noisy joint law of \((V,X)\), so the two
joint laws are at the TV distance of their \(X\)-marginals, at most
\(\tau_n\). Transfer the bounded squared and absolute losses.
\end{proof}

For the set-sector formulation and free-energy gap of the soft code see
Proposition~\ref{ian:hard-sector-soft-instance} in
\path{ian/17_hard_sector_instances.tex}; TV comparison alone is not a
domination statement, and that set-sector argument avoids the issue.

\begin{proposition}[Local orbit families]
\label{ian:phase-label-local-instance}
Let \(H\) and \(k\) be as in Proposition~\ref{ian:hard-sector-local-instances}
(\(n\ge16\), \(\delta_n<1\)). On \(H\), write
\(X_{\rm vis}=\sigma c_A\eta\) with \(c_A\) the nearest character and
\(|\{z:\eta_z=-1\}|<n/32\). Then \(|Q(X)-h(A)|\le1/8\) on \(H\), and
every depth-\(t\) tree \(D\) on all physical spins satisfies
\[
 \E_H|D-Q|
 \ge\frac{1-2^{kt-m/2}}2-\frac18.
 \tag{20.19}\label{ian:phase-label-local-bound}
\]
For the bounded-degree families (sparse three-local, pairwise,
zero-field), the independent-set pinning decomposition (item 5 after
Lemma~\ref{ian:phase-label-close-products}), with \(I\) a set of
visible spins of size at most \(n/8\), is \(\mathcal C_\cup\)-admissible
(using \(s_U\)), its label determines \(h(A)\) on \(H\), and at
least a fraction \(15/17\) of \(\Var_HQ\) is between its labels.
\end{proposition}

\begin{proof}
With \(p_0=\widehat\eta(0)^2\ge225/256\) (proof of
Lemma~\ref{ian:hard-sector-orbit-covariance}),
\(|Q-h(A)|\le1-p_0<1/8\). Conditional on \(H\), \(A\) is uniform and
independent of the side information, and after it is granted each
physical query costs at most \(k\) linear queries to \(A\): this is the
proof of Proposition~\ref{ian:hard-sector-local-instances}, from
Lemma~\ref{res:orbit-lemma} and, for the zero-field family,
Section~\ref{res:zero-field}. The argument for
\eqref{ian:phase-label-clean-bound} with \(kt\) equations and the
triangle inequality give \eqref{ian:phase-label-local-bound}.

For the last assertion, with \(d\) the primal degree of item 5, take
\(I\) among the visible spins with \(n/(d+1)\wedge n/8\le|I|\le n/8\).
Distinct signed characters differ in at least \(n/2\) coordinates, hence
in at least \(3n/8\) outside \(I\), while the visible input has fewer
than \(n/32\) errors. So the label \(X_{I^c}\) determines the nearest
signed character, hence \(h(A)\), on \(H\), and
\(\E_H\Var(Q\mid Z,H)\le\E_H(Q-h(A))^2\le1/64\). With
\(\Var_HQ\ge17/128\) (Proposition~\ref{ian:hard-sector-local-instances}),
the between-label fraction is at least \(1-\frac{1/64}{17/128}=15/17\).
\end{proof}

That pinning label is almost the whole input
(Remark~\ref{ian:phase-label-what-admissibility-excludes}), so this is
not a natural phase label. For the dense three-local family only the
\(\mathcal U\)-decomposition (item 4) is available here; uniform
hereditary simulation of its natural fixed-\(A\) components is not
proved.

\begin{remark}[In which sense each instance is an obstruction]
\label{ian:phase-label-instance-senses}
\emph{Clean code \(U_{C_n}\): only in the \(\widetilde Z\)-sense.} The
natural \(A\)-decomposition is atomic, hence not admissible; the
admissible \(\mathcal U\)-decomposition has negligible between-label
variance \eqref{ian:phase-label-small-between} on the sector \(C_n\), so
no label-only explanation holds on that sector. This is a statement
about the sector \(C_n\) only; on other hard sectors the simplex label
may carry more of \(Q\). Theorem~\ref{ian:phase-label-hard-sector-converse} applies
with \(\widetilde Z=(\xi,\text{one-query transcript})\), which determines
\(Q\).

\emph{Noisy code, label \((A,B)\): label-only.} Take \(H=\Omega_n\),
\(w=1\), and \(\operatorname{Hard}(Q,\nu_{n,\rho};t,a,\rho^4/16)\) with
\(a=3\rho^4/16\), valid when \(2^t/\sqrt n\le\rho^4/16\)
(Proposition~\ref{ian:hard-sector-soft-instance}). By
\eqref{ian:phase-label-noisy-flatness}, \(v\le1/(2n)\le a^2/4096\) once
\(n\ge2^{19}/(9\rho^8)\). Then
\eqref{ian:phase-label-flat-variance}--\eqref{ian:phase-label-flat-covariance}
give \(\Var g_\rho\ge\Var Q-a^2/4096\) and
\(|\Cov(g_\rho,F)|\le\rho^4/16+a/64\) for \(F\in\mathcal F_t\).

\emph{Soft code, label \(V\): label-only.} Take the set \(G\) of
Proposition~\ref{ian:hard-sector-soft-instance}: \(w=\mu(G)\ge1-b_n-\tau_n\),
\(a=\rho^4/8\), \(\gamma=\rho^4/16\). By
\eqref{ian:phase-label-soft-flatness}, \(v\le1/(2n)+\tau_n\to0\), while
\(wa^2/4096\ge(1-b_n-\tau_n)\rho^8/2^{18}\to\rho^8/2^{18}\). Hence
\(v\le wa^2/4096\) for all large \(n\), and
\eqref{ian:phase-label-flat-variance}--\eqref{ian:phase-label-flat-covariance}
hold on \(G\) with \(\sqrt{v/w}\le a/64\).

\emph{Local families.} Equation~\eqref{ian:phase-label-local-bound} is a
label-target bound in the sense of
\eqref{ian:phase-label-pure-target-comparison}, with \(g=h(A)\) and
\(e\le1/8\); the bounded-degree pinning label explains most variance but
is not a phase label.
\end{remark}

\begin{proposition}[Exchangeability and Curie--Weiss are positive cases]
\label{ian:phase-label-positive-instances}
Use the uniform-input simulator \(s_U\), supplied by
Corollary~\ref{cor:resolution-product-alphabet-aa}. Every exchangeable
input law has full AA, with
one polynomial uniform over all exchangeable laws, without recovering
the Hamming-weight label: for \(k=\lfloor n/2\rfloor\), the slices \(k\)
and \(k+1\) are indistinguishable with constant advantage by \(o(n)\)
queries.

For fixed \(\beta>1\), the Curie--Weiss law
\(\mu_{\beta,n}(x)\propto\exp\bigl(\frac{\beta}{2n}(\sum_ix_i)^2\bigr)\)
has a \(\mathcal C_\cup\)-admissible two-phase decomposition with
exchangeable components, and its phase sign is learnable at error
\(\delta'\) with \(O_\beta(\log(1/\delta'))\) queries.
\end{proposition}

\begin{proof}
The exchangeable theorem \ref{exch:heading-2} supplies one tree of cost
polynomial in \((T,\varepsilon^{-1})\) for every slice, hence for every
mixture of slices. For the slices, couple \(X\sim U_{n,k}\) with \(X'\)
obtained by flipping a uniformly chosen zero of \(X\); then
\(X'\sim U_{n,k+1}\), and a depth-\(t\) tree sees a difference only if it
queries the flipped position, which has probability at most
\(t/(n-k)\).

For Curie--Weiss, let \(\mathfrak h(p)=-p\log p-(1-p)\log(1-p)\),
\(\phi(y)=\mathfrak h(\frac{1+y}2)+\frac\beta2y^2\), and let
\(m_\beta\in(0,1)\) solve \(m_\beta=\tanh(\beta m_\beta)\). The mass of
\(\{\sum_ix_i=ny\}\) is \(\binom n{n(1+y)/2}e^{\beta ny^2/2}\) divided by
the normalizer, and
\((n+1)^{-1}e^{n\mathfrak h(j/n)}\le\binom nj\le e^{n\mathfrak h(j/n)}\).
The function \(\phi\) is even, with
\(\phi'(y)=\beta y-\operatorname{artanh}y>0\) on \((0,m_\beta)\), because
\(\tanh(\beta y)>y\) there. So
\(\sup_{|y|\le m_\beta/2}\phi=\phi(m_\beta/2)<\phi(m_\beta)\). An
attainable \(y_n\) lies within \(2/n\) of \(m_\beta\), and \(\phi\) is
Lipschitz, with constant \(c_\beta'\), on
\([m_\beta/2,(1+m_\beta)/2]\); hence, for all large \(n\) and after
enlarging \(C_\beta\) for the others,
\[
 \mu_{\beta,n}\Bigl\{\bigl|\textstyle\sum_iX_i\bigr|<nm_\beta/2\Bigr\}
 \le(n+1)^2e^{2c_\beta'}e^{-n(\phi(m_\beta)-\phi(m_\beta/2))}
 \le C_\beta e^{-c_\beta n}.
\]

Let the label be the sign of \(\sum_ix_i\), with ties assigned to
\(+\). This is a function of \(x\), the two components are exchangeable,
and by the symmetry \(x\mapsto-x\) and the display each has mass at
least \(1/4\) for large \(n\). The conditional field on a spin given the
others is \(\frac\beta n\sum_{j\ne i}x_j\in(-\beta,\beta)\), so every
chain-rule factor of \(\mu_{\beta,n}\) is at most
\((1+e^{-2\beta})^{-1}\), and the component atoms are at most
\(4(1+e^{-2\beta})^{-n}\). With Proposition~\ref{ian:phase-label-union-class}
the decomposition is \(\mathcal C_\cup\)-admissible.

Outside the displayed event, the majority of \(j\) coordinates sampled
with replacement has the correct sign except with probability at most
\(\exp(-jm_\beta^2/8)\), by Hoeffding. Take
\(j=\lceil8m_\beta^{-2}\log(2/\delta')\rceil\). When
\(C_\beta e^{-c_\beta n}>\delta'/2\), \(n=O_\beta(\log(1/\delta'))\), and
reading all spins determines the label. Averaging fixes the sampling
coins.
\end{proof}

The two positive mechanisms differ: the Curie--Weiss phase sign is
learnable, so Theorem~\ref{ian:phase-label-elimination} applies with a
surrogate label, whereas exchangeable simulation uses one predictor
compatible with all slice labels. Neither uses rapid Glauber mixing.
